\documentclass[12pt,a4paper]{article}
\usepackage[top=2.5cm,bottom=2.5cm,left=2.5cm,right=2.5cm,headheight=15pt]{geometry}
\usepackage{amsmath,amssymb}
\usepackage{tikz}
\usetikzlibrary{shapes,arrows,arrows.meta,positioning,calc,backgrounds,fit,matrix}
\usepackage{booktabs,array,longtable,tabularx,multirow}
\usepackage{xcolor}
\usepackage{enumitem}
\usepackage{fancyhdr}
\usepackage{titlesec}
\usepackage{mdframed}
\usepackage{tcolorbox}
\tcbuselibrary{skins,breakable}
\usepackage{hyperref}
\usepackage{forest}

%% ---- Colors ----
\definecolor{folblue}{RGB}{0,82,165}
\definecolor{lightblue}{RGB}{220,235,255}
\definecolor{darkteal}{RGB}{0,105,92}
\definecolor{lightteal}{RGB}{200,230,225}
\definecolor{darkorange}{RGB}{180,70,0}
\definecolor{lightorange}{RGB}{255,235,210}
\definecolor{darkgreen}{RGB}{20,110,20}
\definecolor{lightgreen}{RGB}{200,240,205}
\definecolor{darkred}{RGB}{150,20,20}
\definecolor{lightred}{RGB}{255,210,210}
\definecolor{lightgray}{RGB}{242,242,242}
\definecolor{dpurple}{RGB}{90,0,140}
\definecolor{lpurple}{RGB}{235,215,255}
\definecolor{maroon}{RGB}{128,0,0}

%% ---- tcolorbox styles ----
\newtcolorbox{qbox}[2][]{enhanced,breakable,
  colback=lightorange,colframe=darkorange,arc=4pt,
  title={\textbf{#2}},fonttitle=\bfseries\color{white},
  attach boxed title to top left={yshift=-2mm,xshift=4mm},
  boxed title style={colback=darkorange,arc=3pt},#1}

\newtcolorbox{solbox}[2][]{enhanced,breakable,
  colback=lightblue,colframe=folblue,arc=4pt,
  title={\textbf{#2}},fonttitle=\bfseries\color{white},
  attach boxed title to top left={yshift=-2mm,xshift=4mm},
  boxed title style={colback=folblue,arc=3pt},#1}

\newtcolorbox{keybox}[2][]{enhanced,breakable,
  colback=lightgreen,colframe=darkgreen,arc=4pt,
  title={\textbf{#2}},fonttitle=\bfseries\color{white},
  attach boxed title to top left={yshift=-2mm,xshift=4mm},
  boxed title style={colback=darkgreen,arc=3pt},#1}

\newtcolorbox{algobox}[2][]{enhanced,breakable,
  colback=lightgray,colframe=gray!60,arc=4pt,
  title={\textbf{#2}},fonttitle=\bfseries,
  attach boxed title to top left={yshift=-2mm,xshift=4mm},
  boxed title style={colback=gray!60,arc=3pt},#1}

\newtcolorbox{warnbox}[2][]{enhanced,breakable,
  colback=lightred,colframe=darkred,arc=4pt,
  title={\textbf{#2}},fonttitle=\bfseries\color{white},
  attach boxed title to top left={yshift=-2mm,xshift=4mm},
  boxed title style={colback=darkred,arc=3pt},#1}

%% ---- Page style ----
\pagestyle{fancy}\fancyhf{}
\lhead{\color{folblue}\textbf{FOL \& Resolution --- CSE 713 AI}}
\rhead{\color{folblue}Exam: 10 Jun 2026}
\cfoot{--- \thepage\ ---}
\renewcommand{\headrulewidth}{0.4pt}

%% ---- Section style ----
\titleformat{\section}{\large\bfseries\color{folblue}}{}{0em}{}[\titlerule]
\titleformat{\subsection}{\normalsize\bfseries\color{darkorange}}{}{0em}{}

%% ---- Helpers ----
\newcommand{\ALL}[1]{\forall\,#1\;}
\newcommand{\EX}[1]{\exists\,#1\;}
\newcommand{\ra}{\rightarrow}
\newcommand{\lra}{\leftrightarrow}
\newcommand{\sq}{\square}        % empty clause
\newcommand{\cl}[1]{\{#1\}}      % clause braces
\newcommand{\res}[2]{\textbf{Resolve}\;(#1)\;+\;(#2)}

%% ---- Document ----
\begin{document}

\begin{center}
{\LARGE\bfseries\color{folblue} CSE 713 --- Artificial Intelligence}\\[4pt]
{\large\bfseries FOL, Resolution \& Knowledge Representation}\\[2pt]
{\large Complete Past Paper Solutions: 2020--2024}\\[6pt]
{\normalsize Source: \textit{Propositional Logic.pdf} slides (all 104 pages read)}\\[2pt]
{\normalsize\color{darkred}\textbf{Every question from 2020--2024 included. Zero omissions.}}
\end{center}

\tableofcontents
\newpage

%% ==========================================
\section*{Topic Overview (from Slides)}
%% ==========================================

\begin{keybox}{Topics Covered in This Document (from Propositional Logic.pdf)}
\textbf{Part 1 --- Propositional Logic (PL):}
\begin{itemize}[noitemsep]
  \item Modus Ponens \& valid inference rules
  \item Satisfiability, Validity, Entailment
  \item Clausal form: literals, clauses (disjunctions of literals)
  \item 4-step conversion to clausal form in PL
  \item Resolution rule: $\frac{(x \lor s_1),\;(\neg x \lor s_2)}{s_1 \lor s_2}$
  \item Proof by Refutation (add negated goal $\to$ derive $\square$)
\end{itemize}

\vspace{4pt}\textbf{Part 2 --- First Order Logic (FOL):}
\begin{itemize}[noitemsep]
  \item Limitations of PL (cannot express infinite models)
  \item Terms: Constants, Variables, Functions; Predicates (return T/F)
  \item Quantifiers: $\forall$ (for all), $\exists$ (there exists), duality
  \item FOL sentences: combinations of predicates with $\land,\lor,\neg,\ra,\lra$
  \item Inference rules: Universal Elimination (UE), Existential Elimination/Skolemization (EE), Existential Introduction (EI)
  \item Unification: finding substitution $\theta$ making two atoms identical
  \item 9-step Canonical Form Conversion (CFC) to clause form
  \item Resolution in FOL with Unification
\end{itemize}
\end{keybox}

\begin{algobox}{9-Step Canonical Form Conversion Algorithm (exam-critical)}
\begin{enumerate}[noitemsep,label=\textbf{Step \arabic*.}]
  \item Eliminate biconditionals: $A \lra B \;\equiv\; (A\ra B)\land(B\ra A)$
  \item Eliminate implications: $A \ra B \;\equiv\; \neg A \lor B$
  \item Move $\neg$ inward (De Morgan's, double-negation elimination)
  \item Standardize variables apart (rename so each clause has unique vars)
  \item Move all quantifiers to the left (prenex normal form)
  \item Eliminate existential quantifiers (\textbf{Skolemization}):
    \begin{itemize}[noitemsep]
      \item $\exists y$ outside any $\forall$: replace $y$ with a Skolem \textbf{constant} $c$
      \item $\forall x\;\exists y$: replace $y$ with a Skolem \textbf{function} $f(x)$
    \end{itemize}
  \item Drop all $\forall$ (variables implicitly universally quantified)
  \item Distribute $\land$ over $\lor$ (convert to CNF)
  \item Write each conjunct as a separate clause (set of literals)
\end{enumerate}
\end{algobox}

\newpage

%% ==========================================
\section{2020 --- Group-A Q4: Marcus/Pompeii (approx.\ 9 marks)}
%% ==========================================

\begin{qbox}{Question --- 2020 Group-A Q4}
Consider the following facts:
\begin{enumerate}[noitemsep,label=(\roman*)]
  \item Marcus was a man.
  \item Marcus was a Pompeian.
  \item Marcus was born in 40 A.D.
  \item All men are mortal.
  \item All Pompeians died when the volcano erupted in 79 A.D.
  \item No mortal lives longer than 150 years.
  \item It is now 2021.
  \item Alive means not dead.
  \item If someone dies, then he is dead at all later times.
\end{enumerate}
\begin{enumerate}[label=(\alph*),noitemsep]
  \item Translate these facts into well-formed formulas (WFFs) in predicate logic.
  \item Answer the question ``Is Marcus alive now?'' using backward reasoning.
  \item Convert the formulas into clause form.
  \item Prove that ``Marcus is not alive now'' using resolution.
\end{enumerate}
\end{qbox}

\subsection{Part (a): WFF Translation}

\begin{solbox}{WFF Translation --- 2020 Q4(a)}
\begin{align*}
  F_1&: \quad \mathrm{man(Marcus)} \\[2pt]
  F_2&: \quad \mathrm{pompeian(Marcus)} \\[2pt]
  F_3&: \quad \mathrm{born(Marcus,\;40)} \\[2pt]
  F_4&: \quad \ALL{x}\;\mathrm{man}(x) \ra \mathrm{mortal}(x) \\[2pt]
  F_5&: \quad \ALL{x}\;\mathrm{pompeian}(x) \ra \mathrm{dies}(x,\;79) \\[2pt]
  F_6&: \quad \ALL{x}\ALL{b}\ALL{t}\;
         \mathrm{mortal}(x) \land \mathrm{born}(x,b) \land \mathrm{gt}(t,\;\mathrm{plus}(b,150))
         \ra \mathrm{dead}(x,t) \\[2pt]
  F_7&: \quad \mathrm{now} = 2021 \\[2pt]
  F_8&: \quad \ALL{x}\;\mathrm{alive}(x) \lra \neg\,\mathrm{dead}(x,\;\mathrm{now}) \\[2pt]
  F_9&: \quad \ALL{x}\ALL{t_1}\ALL{t_2}\;
         \mathrm{dies}(x,t_1) \land \mathrm{gt}(t_2,t_1) \ra \mathrm{dead}(x,t_2)
\end{align*}
\end{solbox}

\subsection{Part (b): Backward Reasoning --- Is Marcus Alive?}

\begin{solbox}{Backward Reasoning --- 2020 Q4(b)}
\textbf{Goal:} Determine \textit{alive}(Marcus).\quad
Start from the goal and work backwards through the rules.
\begin{enumerate}[noitemsep,label=\textbf{Step \arabic*.}]
  \item From $F_8$: $\mathrm{alive(Marcus)} \lra \neg\,\mathrm{dead(Marcus, now)}$.\\
        $\Rightarrow$ Need to determine whether $\mathrm{dead(Marcus, 2021)}$ is true.
  \item To prove $\mathrm{dead(Marcus,\;2021)}$, use $F_9$ (with $x/\mathrm{Marcus},\;t_2/2021$):\\
        Need: $\mathrm{dies(Marcus,\;t_1)} \land \mathrm{gt}(2021,\;t_1)$ for some $t_1$.
  \item From $F_5$ (with $x/\mathrm{Marcus}$): $\mathrm{pompeian(Marcus)} \ra \mathrm{dies(Marcus,\;79)}$.\\
        From $F_2$: $\mathrm{pompeian(Marcus)}$ is true.\quad
        $\therefore\;\mathrm{dies(Marcus,\;79)}$. \quad Set $t_1 = 79$.
  \item $\mathrm{gt}(2021,\;79) = \mathrm{TRUE}$ \quad (arithmetic fact).
  \item $\therefore\;\mathrm{dead(Marcus,\;2021)}$ (i.e., $\mathrm{dead(Marcus,\;now)}$) is TRUE.
  \item From $F_8$: $\mathrm{alive(Marcus)} \lra \neg\,\mathrm{dead(Marcus,\;now)}$\\
        Since $\mathrm{dead(Marcus,\;now)}$ is TRUE, $\;\mathrm{alive(Marcus)} = \mathbf{FALSE}$.
\end{enumerate}
\textbf{Answer: Marcus is NOT alive now.}
\end{solbox}

\subsection{Part (c): Clause Form Conversion}

\begin{solbox}{Clause Form --- 2020 Q4(c)}
Apply the 9-step CFC algorithm:

\medskip
\begin{tabular}{@{}l@{\quad}l@{\quad}l@{}}
  \toprule
  \textbf{Formula} & \textbf{Steps Applied} & \textbf{Clause(s)} \\
  \midrule
  $F_1$: man(Marcus) & ground atom & $C_1:\cl{\mathrm{man(Marcus)}}$ \\[3pt]
  $F_2$: pompeian(Marcus) & ground atom & $C_2:\cl{\mathrm{pompeian(Marcus)}}$ \\[3pt]
  $F_3$: born(Marcus,40) & ground atom & $C_3:\cl{\mathrm{born(Marcus,40)}}$ \\[3pt]
  $F_4$: $\ALL{x}\mathrm{man}(x)\ra\mathrm{mortal}(x)$ & Steps 2,7 & $C_4:\cl{\neg\mathrm{man}(x),\;\mathrm{mortal}(x)}$ \\[3pt]
  $F_5$: $\ALL{x}\mathrm{pompeian}(x)\ra\mathrm{dies}(x,79)$ & Steps 2,7 & $C_5:\cl{\neg\mathrm{pompeian}(x),\;\mathrm{dies}(x,79)}$ \\[3pt]
  $F_6$: mortal+born+gt$\ra$dead & Steps 2,7 & $C_6:\cl{\neg\mathrm{mortal}(x),\neg\mathrm{born}(x,b),\neg\mathrm{gt}(t,\mathrm{plus}(b,150)),\mathrm{dead}(x,t)}$ \\[3pt]
  $F_7$: now=2021 & fact & $C_7:\cl{\mathrm{now=2021}}$ \;(arithmetic) \\[3pt]
  $F_8$: $\ALL{x}\mathrm{alive}(x)\lra\neg\mathrm{dead}(x,\mathrm{now})$ & Steps 1,2,7,9 &
    $C_8:\cl{\neg\mathrm{alive}(x),\;\neg\mathrm{dead}(x,\mathrm{now})}$ \\
    & & $C_9:\cl{\mathrm{dead}(x,\mathrm{now}),\;\mathrm{alive}(x)}$ \\[3pt]
  $F_9$: dies$\land$gt$\ra$dead & Steps 2,7 & $C_{10}:\cl{\neg\mathrm{dies}(x,t_1),\;\neg\mathrm{gt}(t_2,t_1),\;\mathrm{dead}(x,t_2)}$ \\
  \bottomrule
\end{tabular}
\end{solbox}

\subsection{Part (d): Resolution Proof --- Marcus is NOT Alive Now}

\begin{solbox}{Resolution Proof by Refutation --- 2020 Q4(d)}
\textbf{Strategy:} Negate the goal. Add $C_{\neg}:\cl{\mathrm{alive(Marcus)}}$ and derive the empty clause $\square$.

\medskip
\begin{tabular}{@{}l@{\quad}l@{\quad}l@{\quad}l@{}}
  \toprule
  \textbf{Step} & \textbf{Clauses Used} & \textbf{Substitution} & \textbf{Result} \\
  \midrule
  1 & $C_1\;\&\;C_4$ & $\{x/\mathrm{Marcus}\}$ & $C_{11}:\cl{\mathrm{mortal(Marcus)}}$ \\[3pt]
  2 & $C_2\;\&\;C_5$ & $\{x/\mathrm{Marcus}\}$ & $C_{12}:\cl{\mathrm{dies(Marcus,79)}}$ \\[3pt]
  3 & $C_{12}\;\&\;C_{10}$ & $\{x/\mathrm{Marcus},t_1/79\}$ & $C_{13}:\cl{\neg\mathrm{gt}(t_2,79),\;\mathrm{dead(Marcus},t_2)}$ \\[3pt]
  4 & $C_{13}$ + gt(2021,79) true & $\{t_2/2021\}$ & $C_{14}:\cl{\mathrm{dead(Marcus,2021)}}$ \\
    & & (now=2021) & $\equiv\cl{\mathrm{dead(Marcus,now)}}$ \\[3pt]
  5 & $C_{\neg}\;\&\;C_8$ & $\{x/\mathrm{Marcus}\}$ & $C_{15}:\cl{\neg\mathrm{dead(Marcus,now)}}$ \\[3pt]
  6 & $C_{14}\;\&\;C_{15}$ & --- & $\boxed{\square}$ \textbf{(Contradiction!)} \\
  \bottomrule
\end{tabular}

\medskip
\textbf{Conclusion:} Adding alive(Marcus) to the KB leads to contradiction.\\
$\therefore\;\mathrm{alive(Marcus)}$ is FALSE, i.e., \textbf{Marcus is NOT alive now.} \checkmark
\end{solbox}

\newpage
%% ==========================================
\section{2021 --- Q4: PL Reasoning (approx.\ 9.25 marks)}
%% ==========================================

\begin{qbox}{Question --- 2021 Q4}
\begin{enumerate}[label=(\alph*),noitemsep]
  \item The figure shows a Wumpus World game. The agent starts from [1,1]. Briefly discuss with diagrams how the agent generates new knowledge through perceiving the environment and choosing its next move. \hfill\textit{(3 marks)}
  \item Differentiate between proof and theorem in one point. \hfill\textit{(0.75 marks)}
  \item Determine whether the following sentence is (i) Satisfiable, (ii) Contradictory, or (iii) Valid:\\
        \[(P \lor Q) \land (P \lor \neg Q) \lor P\] \hfill\textit{(2 marks)}
  \item Given knowledge base $R_1$--$R_5$ below, show that $P_{1,2}$ is false using resolution: \hfill\textit{(3.5 marks)}\\
        $R_1: \neg P_{1,1}$\quad $R_2: B_{1,1} \lra (P_{2,1} \lor P_{1,2})$\quad $R_3: B_{2,1} \lra (P_{1,1} \lor P_{2,2} \lor P_{3,1})$\\
        $R_4: \neg B_{1,1}$\quad $R_5: B_{2,1}$
\end{enumerate}
\end{qbox}

\subsection{Part (a): Wumpus World Knowledge Generation (3 marks)}

\begin{solbox}{Wumpus World --- 2021 Q4(a)}
The Wumpus World uses \textbf{Propositional Logic (PL)} to reason about hidden dangers. The agent builds a KB incrementally from perceptions.

\medskip
\textbf{Key Propositions:}
\begin{itemize}[noitemsep]
  \item $B_{x,y}$: Breeze is felt at cell $[x,y]$ \quad $P_{x,y}$: Pit exists at $[x,y]$
  \item $S_{x,y}$: Stench is felt at $[x,y]$ \quad $W_{x,y}$: Wumpus exists at $[x,y]$
\end{itemize}

\textbf{Initial KB rules (KB axioms):}
\begin{align*}
  &R_1: \neg P_{1,1} \quad\text{(no pit at start)}\\
  &R_2: B_{1,1} \lra (P_{2,1} \lor P_{1,2})\\
  &R_4: \neg B_{1,1} \quad\text{(agent perceives no breeze at [1,1])}
\end{align*}

\textbf{Inference chain:}
\begin{enumerate}[noitemsep,label=\textbf{\arabic*.}]
  \item Agent at [1,1]: perceives $\neg B_{1,1}$ (no breeze), $\neg S_{1,1}$ (no stench).
  \item From $\neg B_{1,1}$ + $R_2$ (biconditional left-to-right negated):\\
        $\neg B_{1,1} \ra \neg(P_{2,1} \lor P_{1,2}) \equiv \neg P_{2,1} \land \neg P_{1,2}$\\
        $\therefore$ Cells [2,1] and [1,2] are \textbf{PIT-FREE} (safe to enter).
  \item Agent moves to [2,1]: perceives $B_{2,1}$ (breeze felt). From $R_5: B_{2,1}$ is added to KB.
  \item From $R_3$: $B_{2,1} \lra (P_{1,1} \lor P_{2,2} \lor P_{3,1})$\\
        $B_{2,1}$ is TRUE $\Rightarrow$ at least one of $\{P_{1,1}, P_{2,2}, P_{3,1}\}$ is true.\\
        $\neg P_{1,1}$ (from $R_1$) is known. $\neg P_{2,2}$ can be inferred (no path explored).\\
        $\therefore\;P_{3,1}$ is TRUE --- \textbf{Pit at [3,1]}!
  \item Agent \textbf{avoids [3,1]}, backtracks, and continues to [1,2] (safe from Step 2).
\end{enumerate}

Each perception updates the KB; logical inference (resolution) derives new facts about safe/dangerous cells without direct observation.
\end{solbox}

\subsection{Part (b): Proof vs.\ Theorem (0.75 marks)}

\begin{solbox}{Proof vs.\ Theorem --- 2021 Q4(b)}
\begin{itemize}[noitemsep]
  \item A \textbf{theorem} is a statement (proposition) that has been established as true within a formal system.
  \item A \textbf{proof} is the sequence of logical steps (inference rules, axioms) that \textit{establishes} the truth of a theorem.
\end{itemize}
\textit{In one point:} A theorem is the \textbf{conclusion}; a proof is the \textbf{derivation} that justifies it.
\end{solbox}

\subsection{Part (c): Satisfiability Analysis (2 marks)}

\begin{solbox}{Satisfiability of $(P \lor Q)\land(P \lor \neg Q)\lor P$ --- 2021 Q4(c)}
\textbf{Step 1: Apply operator precedence} ($\neg > \land > \lor$):
\[(P \lor Q)\land(P \lor \neg Q)\lor P \;=\; \bigl[(P \lor Q)\land(P \lor \neg Q)\bigr] \lor P\]

\textbf{Step 2: Simplify the bracketed expression} using Distributivity:
\[(P \lor Q)\land(P \lor \neg Q) \;=\; P \lor (Q \land \neg Q) \;=\; P \lor \mathrm{F} \;=\; P\]

\textbf{Step 3: Final expression:}
\[P \lor P \;=\; P\]

\textbf{Step 4: Truth table check:}
\begin{center}
\begin{tabular}{c|c}
  \toprule
  $P$ & Expression value\\
  \midrule
  T & T \\
  F & F \\
  \bottomrule
\end{tabular}
\end{center}

\textbf{Conclusion:}
\begin{itemize}[noitemsep]
  \item \textbf{Satisfiable} --- TRUE when $P = \mathrm{T}$. \checkmark
  \item NOT Contradictory (not always false)
  \item NOT Valid (not always true; false when $P = \mathrm{F}$)
\end{itemize}
\end{solbox}

\subsection{Part (d): Resolution Proof --- $P_{1,2}$ is False (3.5 marks)}

\begin{solbox}{KB Resolution --- Show $P_{1,2}$ is FALSE --- 2021 Q4(d)}
\textbf{Given KB:}
\begin{align*}
  R_1&: \neg P_{1,1} \\
  R_2&: B_{1,1} \lra (P_{2,1} \lor P_{1,2})\\
  R_3&: B_{2,1} \lra (P_{1,1} \lor P_{2,2} \lor P_{3,1})\\
  R_4&: \neg B_{1,1}\\
  R_5&: B_{2,1}
\end{align*}

\textbf{Convert $R_2$ to clause form} (Steps 1--2 of CFC):
\begin{align*}
  R_2 \equiv &\underbrace{(B_{1,1} \ra P_{2,1} \lor P_{1,2})}_{C_a:\;\neg B_{1,1} \lor P_{2,1} \lor P_{1,2}}
  \land
  \underbrace{(P_{2,1} \lor P_{1,2} \ra B_{1,1})}_{C_b:\;\neg P_{2,1} \lor B_{1,1}\;\text{and}\;C_c:\;\neg P_{1,2} \lor B_{1,1}}
\end{align*}

\textbf{Goal:} Prove $\neg P_{1,2}$ (i.e., $P_{1,2}$ is false).\\
\textbf{Method:} Refutation --- add $\cl{P_{1,2}}$ (negation of goal) to KB, derive $\square$.

\medskip
\begin{tabular}{@{}l@{\quad}l@{\quad}l@{}}
  \toprule
  \textbf{Step} & \textbf{Clauses Resolved} & \textbf{Result} \\
  \midrule
  1 & $\cl{P_{1,2}}$ + $C_c:\cl{\neg P_{1,2} \lor B_{1,1}}$ & $\cl{B_{1,1}}$ \\
  2 & $\cl{B_{1,1}}$ + $R_4:\cl{\neg B_{1,1}}$ & $\boxed{\square}$ \textbf{(Contradiction!)}\\
  \bottomrule
\end{tabular}

\medskip
Since assuming $P_{1,2}$ leads to contradiction, \textbf{$P_{1,2}$ is FALSE}. \checkmark

\medskip
\textbf{Intuition:} $R_4$ says no breeze at [1,1]. $R_2$ then tells us no pit at [2,1] or [1,2].
\end{solbox}

\newpage
%% ==========================================
\section{2021 --- Q5(a,b): FOL Generalization + Perfect Square (approx.\ 8.75 marks)}
%% ==========================================

\begin{qbox}{Question --- 2021 Q5(a) \& Q5(b)}
\begin{enumerate}[label=(\alph*),noitemsep]
  \item Briefly explain why FOL is considered as the generalization of PL. \hfill\textit{(1 mark)}
  \item Consider the following sentences: \hfill\textit{(7.75 marks)}
    \begin{enumerate}[label=(\roman*),noitemsep]
      \item If a perfect square is divisible by a prime $P$, then it is also divisible by the square of $P$.
      \item Every perfect square is divisible by some prime.
      \item 36 is a perfect square.
      \item Does there exist a prime $q$ such that the square of $q$ divides 36?
    \end{enumerate}
    Prove via the resolution method and unification (where needed) that ``Does there exist a prime $q$ such that the square of $q$ divides 36?'' More precisely:
    \begin{itemize}[noitemsep]
      \item[(i)] First, translate the sentences above into FOL sentences.
      \item[(ii)] Convert to Clause Forms.
      \item[(iii)] Apply the resolution method with unification (where needed) to prove the goal sentence.
    \end{itemize}
\end{enumerate}
\end{qbox}

\subsection{Part (a): FOL as Generalization of PL (1 mark)}

\begin{solbox}{FOL vs.\ PL --- 2021 Q5(a)}
\textbf{PL} can only express ground propositions (true/false facts about specific named objects). It \textbf{cannot} handle statements about all members of a set without listing them individually.

\textbf{FOL extends PL} by adding:
\begin{enumerate}[noitemsep]
  \item \textbf{Variables} ranging over objects (not just truth values)
  \item \textbf{Universal quantifier} $\forall$: ``for all objects of a given type''
  \item \textbf{Existential quantifier} $\exists$: ``there exists at least one object''
  \item \textbf{Predicates} (relations on objects) and \textbf{Functions} (maps objects to objects)
\end{enumerate}
Example: ``All men are mortal'' requires infinitely many PL sentences ($\mathrm{mortal(Marcus)}$, $\mathrm{mortal(Brutus)}$, \ldots) but only ONE FOL sentence: $\ALL{x}\,\mathrm{man}(x) \ra \mathrm{mortal}(x)$.
\end{solbox}

\subsection{Part (b-i): FOL Translation (1.75 marks)}

\begin{solbox}{FOL Translation --- 2021 Q5(b-i)}
\textbf{Predicates declared:}
\begin{itemize}[noitemsep]
  \item $\mathrm{PerfSq}(x)$: $x$ is a perfect square
  \item $\mathrm{Prime}(p)$: $p$ is a prime number
  \item $\mathrm{Divides}(a,x)$: $a$ divides $x$
  \item $\mathrm{square}(p)$: function returning $p^2$
\end{itemize}

\begin{align}
  F_1&: \quad \ALL{x}\ALL{p}\;\mathrm{PerfSq}(x) \land \mathrm{Prime}(p) \land \mathrm{Divides}(p,x)
          \ra \mathrm{Divides}(\mathrm{square}(p),\;x) \label{f1}\\[4pt]
  F_2&: \quad \ALL{x}\;\mathrm{PerfSq}(x) \ra \EX{p}\;\mathrm{Prime}(p) \land \mathrm{Divides}(p,x) \label{f2}\\[4pt]
  F_3&: \quad \mathrm{PerfSq}(36) \label{f3}\\[4pt]
  \text{Goal}&: \quad \EX{q}\;\mathrm{Prime}(q) \land \mathrm{Divides}(\mathrm{square}(q),\;36) \label{goal}
\end{align}
\end{solbox}

\subsection{Part (b-ii): Canonical Form / Clause Form (2.5 marks)}

\begin{solbox}{Clause Form Conversion --- 2021 Q5(b-ii)}

\textbf{Clause from $F_1$} (Steps 2, 7):
\begin{align*}
  F_1 &\ra \neg\mathrm{PerfSq}(x) \lor \neg\mathrm{Prime}(p) \lor \neg\mathrm{Divides}(p,x) \lor \mathrm{Divides}(\mathrm{square}(p),x)\\
  &\Rightarrow \quad C_1: \;\cl{\neg\mathrm{PerfSq}(x),\;\neg\mathrm{Prime}(p),\;\neg\mathrm{Divides}(p,x),\;\mathrm{Divides}(\mathrm{square}(p),x)}
\end{align*}

\textbf{Clause from $F_2$} (Steps 2, 6 Skolemization, 7):

$F_2$: $\ALL{x}\,\mathrm{PerfSq}(x) \ra \EX{p}\,\mathrm{Prime}(p)\land\mathrm{Divides}(p,x)$

The $\exists p$ is inside $\forall x$, so replace $p$ with Skolem function $g(x)$:
\begin{align*}
  &\ALL{x}\;\mathrm{PerfSq}(x) \ra \mathrm{Prime}(g(x)) \land \mathrm{Divides}(g(x),x)\\
  \text{Split:}\quad
  &C_{2a}: \;\cl{\neg\mathrm{PerfSq}(x),\;\mathrm{Prime}(g(x))}\\
  &C_{2b}: \;\cl{\neg\mathrm{PerfSq}(x),\;\mathrm{Divides}(g(x),x)}
\end{align*}

\textbf{Clause from $F_3$:}
\[C_3: \;\cl{\mathrm{PerfSq}(36)}\]

\textbf{Negated Goal} (negate $F_{\text{goal}}$ for refutation):
\[\neg\EX{q}\,\mathrm{Prime}(q)\land\mathrm{Divides}(\mathrm{square}(q),36)
  \;\equiv\; \ALL{q}\,\neg\mathrm{Prime}(q) \lor \neg\mathrm{Divides}(\mathrm{square}(q),36)\]
\[C_4: \;\cl{\neg\mathrm{Prime}(q),\;\neg\mathrm{Divides}(\mathrm{square}(q),36)}\]
\end{solbox}

\subsection{Part (b-iii): Resolution Proof with Unification (3.5 marks)}

\begin{solbox}{Resolution Proof --- 2021 Q5(b-iii)}
\textbf{Clause set:}
\quad $C_1$, $C_{2a}$, $C_{2b}$, $C_3$ (KB) $+$ $C_4$ (negated goal)

\medskip
\begin{tabular}{@{}l@{\;\;}l@{\;\;}l@{\;\;}l@{}}
  \toprule
  \textbf{Step} & \textbf{Resolved} & \textbf{Unifier} & \textbf{Result} \\
  \midrule
  1 & $C_3 + C_{2a}$ & $\{x/36\}$ & $C_5: \cl{\mathrm{Prime}(g(36))}$ \\[3pt]
  2 & $C_3 + C_{2b}$ & $\{x/36\}$ & $C_6: \cl{\mathrm{Divides}(g(36),\;36)}$ \\[3pt]
  3 & $C_3 + C_1$ & $\{x/36\}$ & $C_7: \cl{\neg\mathrm{Prime}(p),\;\neg\mathrm{Divides}(p,36),\;\mathrm{Divides}(\mathrm{square}(p),36)}$ \\[3pt]
  4 & $C_5 + C_7$ & $\{p/g(36)\}$ & $C_8: \cl{\neg\mathrm{Divides}(g(36),36),\;\mathrm{Divides}(\mathrm{square}(g(36)),36)}$ \\[3pt]
  5 & $C_6 + C_8$ & --- & $C_9: \cl{\mathrm{Divides}(\mathrm{square}(g(36)),\;36)}$ \\[3pt]
  6 & $C_5 + C_4$ & $\{q/g(36)\}$ & $C_{10}: \cl{\neg\mathrm{Divides}(\mathrm{square}(g(36)),\;36)}$ \\[3pt]
  7 & $C_9 + C_{10}$ & --- & $\boxed{\square}$ \textbf{(Contradiction!)} \\
  \bottomrule
\end{tabular}

\medskip
\textbf{Conclusion:} Negating the goal leads to contradiction. Therefore:\\
$\exists\,q$ such that $\mathrm{Prime}(q) \land \mathrm{Divides}(q^2, 36)$ is \textbf{TRUE}. \checkmark

\medskip
\textit{Note:} $g(36)$ is the Skolem function instantiated at 36. It represents the specific prime (e.g., 2 or 3) that divides 36. In reality, $2^2 = 4$ divides 36, and $3^2 = 9$ divides 36.
\end{solbox}

\newpage
%% ==========================================
\section{2022 --- A-2: Resolution + Marcus/Pompeii (7 marks)}
%% ==========================================

\begin{qbox}{Question --- 2022 A-2}
\begin{enumerate}[label=(\alph*),noitemsep]
  \item What is resolution? Write down the algorithm for the conversion of well-formed formulas (wff) in predicate logic to clause form. \hfill\textit{(2.5 marks)}
  \item Consider the following facts (F-1 to F-9) about Marcus and Pompeii. ``It is now 2023.''\\
    (i) Translate these facts into wff in predicate logic. \hfill\textit{(1 mark)}\\
    (ii) Answer the question ``Is Marcus alive now?'' using backward reasoning. \hfill\textit{(1 mark)}\\
    (iii) Convert the formula into clause form. \hfill\textit{(1 mark)}\\
    (iv) Prove that ``Marcus is not alive now'' using resolution. \hfill\textit{(2 marks)}
  \item What do you mean by Knowledge Representation and Mapping? What are the roles of Knowledge Representation? \hfill\textit{(1.5 marks)}
\end{enumerate}
\end{qbox}

\subsection{Part (a): Resolution + CFC Algorithm (2.5 marks)}

\begin{solbox}{Resolution + CFC Algorithm --- 2022 A-2(a)}
\textbf{Resolution} is a refutation-complete inference procedure that operates on clauses.\\
Given two clauses containing complementary literals ($L$ and $\neg L$), resolution derives a new clause (the \textit{resolvent}) with those literals removed:
\[\frac{(L \lor \alpha_1 \lor \cdots \lor \alpha_n)\quad(\neg L \lor \beta_1 \lor \cdots \lor \beta_m)}
       {(\alpha_1 \lor \cdots \lor \alpha_n \lor \beta_1 \lor \cdots \lor \beta_m)}\]
In FOL, we use \textbf{unification} to find a substitution $\theta$ making $L\theta = L'\theta$ before resolving.

\textbf{Algorithm for wff $\to$ Clause Form (9 steps):}
\begin{enumerate}[noitemsep,label=\textbf{Step \arabic*.}]
  \item \textbf{Eliminate biconditionals}: $A \lra B \equiv (A \ra B) \land (B \ra A)$
  \item \textbf{Eliminate implications}: $A \ra B \equiv \neg A \lor B$
  \item \textbf{Move $\neg$ inward}: Apply De Morgan's laws; eliminate double negations
  \item \textbf{Standardize variables}: Rename vars so each clause has distinct variable names
  \item \textbf{Prenex normal form}: Move all quantifiers to the left
  \item \textbf{Skolemization} (eliminate $\exists$):
    \begin{itemize}[noitemsep]
      \item $\exists y$ outside all $\forall$: replace $y$ with Skolem constant $c$
      \item $\forall x_1\cdots x_n\,\exists y$: replace $y$ with Skolem function $f(x_1,\ldots,x_n)$
    \end{itemize}
  \item \textbf{Drop all $\forall$}
  \item \textbf{Distribute $\land$ over $\lor$} (convert to CNF)
  \item \textbf{Write as set of clauses} (each conjunct is a clause = set of literals)
\end{enumerate}
\end{solbox}

\subsection{Part (b): Marcus/Pompeii 2022 (``now = 2023'') (5 marks)}

\begin{solbox}{WFF Translation --- 2022 A-2(b-i)}
\textit{Same as 2020 but with $\mathrm{now} = 2023$:}
\begin{align*}
  F_1&: \mathrm{man(Marcus)}\quad F_2: \mathrm{pompeian(Marcus)}\quad F_3: \mathrm{born(Marcus,40)}\\
  F_4&: \ALL{x}\,\mathrm{man}(x)\ra\mathrm{mortal}(x)\\
  F_5&: \ALL{x}\,\mathrm{pompeian}(x)\ra\mathrm{dies}(x,79)\\
  F_6&: \ALL{x}\ALL{b}\ALL{t}\,\mathrm{mortal}(x)\land\mathrm{born}(x,b)\land\mathrm{gt}(t,\mathrm{plus}(b,150))\ra\mathrm{dead}(x,t)\\
  F_7&: \mathrm{now}=2023\\
  F_8&: \ALL{x}\,\mathrm{alive}(x)\lra\neg\mathrm{dead}(x,\mathrm{now})\\
  F_9&: \ALL{x}\ALL{t_1}\ALL{t_2}\,\mathrm{dies}(x,t_1)\land\mathrm{gt}(t_2,t_1)\ra\mathrm{dead}(x,t_2)
\end{align*}
\end{solbox}

\begin{solbox}{Backward Reasoning --- 2022 A-2(b-ii)}
Identical chain to 2020, but $\mathrm{now}=2023$. Since $\mathrm{gt}(2023,79)$ is TRUE, the chain holds:
\[\mathrm{pompeian(Marcus)}\xrightarrow{F_5}\mathrm{dies(Marcus,79)}\xrightarrow{F_9,\;\mathrm{gt}(2023,79)}\mathrm{dead(Marcus,2023)}\xrightarrow{F_8}\neg\mathrm{alive(Marcus)}\]
\textbf{Answer: Marcus is NOT alive in 2023.}
\end{solbox}

\begin{solbox}{Clause Form --- 2022 A-2(b-iii)}
Clause set is identical to 2020 (Section 1, Part c), with $\mathrm{now}=2023$. Key clauses:
\begin{align*}
  C_4&: \cl{\neg\mathrm{man}(x),\;\mathrm{mortal}(x)}\\
  C_5&: \cl{\neg\mathrm{pompeian}(x),\;\mathrm{dies}(x,79)}\\
  C_8&: \cl{\neg\mathrm{alive}(x),\;\neg\mathrm{dead}(x,\mathrm{now})}\\
  C_{10}&: \cl{\neg\mathrm{dies}(x,t_1),\;\neg\mathrm{gt}(t_2,t_1),\;\mathrm{dead}(x,t_2)}
\end{align*}
\end{solbox}

\begin{solbox}{Resolution Proof --- 2022 A-2(b-iv)}
\textit{Same steps as 2020 Section 1 Part (d), with now=2023.}\\
Refutation: Add $\cl{\mathrm{alive(Marcus)}}$, derive $\square$ in 6 steps. (See Section 1 for full trace.)\\
\textbf{Conclusion: Marcus is NOT alive now.} \checkmark
\end{solbox}

\subsection{Part (c): Knowledge Representation and Mapping (1.5 marks)}

\begin{solbox}{Knowledge Representation --- 2022 A-2(c)}
\textbf{Knowledge Representation (KR)} is the field concerned with encoding world knowledge in a formal structure that AI systems can use to reason, solve problems, and make decisions.

\textbf{Mapping} refers to the correspondence between real-world entities/relationships and their formal symbols in the KR system. For example, the fact ``Marcus is a man'' maps to the atom $\mathrm{man(Marcus)}$.

\textbf{Roles of Knowledge Representation:}
\begin{enumerate}[noitemsep]
  \item \textbf{Surrogate}: Substitutes for real-world entities, enabling reasoning without physical interaction.
  \item \textbf{Ontological commitment}: Defines what exists in the domain (the ``things'' the system knows about).
  \item \textbf{Inferential foundation}: Supports derivation of new facts from existing ones via inference rules.
  \item \textbf{Medium of expression}: A formal language for humans to communicate domain knowledge to machines.
  \item \textbf{Efficient computation}: Structured representations (clauses, networks) support efficient reasoning algorithms.
\end{enumerate}
\end{solbox}

\newpage
%% ==========================================
\section{2023 --- Q4: FOL Translation + Drone (9 marks)}
%% ==========================================

\begin{qbox}{Question --- 2023 Q4}
\begin{enumerate}[label=(\alph*),noitemsep]
  \item Define the following actions of a self-driving car using First-Order Logic (FOL). Declare the predicates first. \hfill\textit{(5 marks)}
    \begin{enumerate}[label=(\roman*),noitemsep]
      \item The car can move to a destination if the destination is reachable and there are no obstacles on the road.
      \item A self-driving car can proceed through an intersection if the traffic light is green and there are no pedestrians crossing.
      \item How would you represent in FOL that a self-driving car should stop if there is an obstacle on the road?
      \item A self-driving car should reduce its speed when entering a school zone if the school zone is active (e.g., during certain hours of the day).
      \item Represent in FOL the rule that a self-driving car can only change lanes if the adjacent lane is clear and the car is moving.
    \end{enumerate}
  \item A drone is responsible for delivering a package to a destination. Drone starts at $L_0$, air route $R_1$ connects $L_0$ to $L_1$, $R_2$ connects $L_1$ to $L_2$ (\textbf{obstacle on $R_2$}), and $R_3$ connects $L_2$ to $L_p$ (destination). Define the predicates, knowledge base, and rules in FOL, then use forward chaining to determine if the drone can reach the destination. \hfill\textit{(4 marks)}
\end{enumerate}
\end{qbox}

\subsection{Part (a): Self-Driving Car FOL (5 marks)}

\begin{solbox}{Self-Driving Car FOL --- 2023 Q4(a)}
\textbf{Predicate declarations:}
\begin{itemize}[noitemsep]
  \item $\mathrm{Car}(c)$: $c$ is a self-driving car
  \item $\mathrm{Reachable}(c,d)$: destination $d$ is reachable by car $c$
  \item $\mathrm{Obstacle}(d)$: there is an obstacle at/on road to $d$
  \item $\mathrm{Moves}(c,d)$: car $c$ moves to destination $d$
  \item $\mathrm{AtIntersection}(c,i)$: car $c$ is at intersection $i$
  \item $\mathrm{GreenLight}(i)$: traffic light at $i$ is green
  \item $\mathrm{PedestrianCrossing}(i)$: pedestrian is crossing at $i$
  \item $\mathrm{Proceeds}(c,i)$: car $c$ proceeds through intersection $i$
  \item $\mathrm{ObstacleOnRoad}(c)$: obstacle on road ahead of car $c$
  \item $\mathrm{Stops}(c)$: car $c$ stops
  \item $\mathrm{InSchoolZone}(c,z)$: car $c$ is in school zone $z$
  \item $\mathrm{ActiveHours}(z)$: school zone $z$ is currently active
  \item $\mathrm{ReduceSpeed}(c)$: car $c$ reduces speed
  \item $\mathrm{AdjacentLaneClear}(c,l)$: lane $l$ adjacent to car $c$'s lane is clear
  \item $\mathrm{Moving}(c)$: car $c$ is moving
  \item $\mathrm{ChangeLane}(c,l)$: car $c$ changes to lane $l$
\end{itemize}

\medskip\textbf{(i)} The car can move to a destination if reachable and no obstacles:
\[\ALL{c}\ALL{d}\;\mathrm{Car}(c) \land \mathrm{Reachable}(c,d) \land \neg\mathrm{Obstacle}(d) \ra \mathrm{Moves}(c,d)\]

\textbf{(ii)} Proceed through intersection if green light and no pedestrians:
\[\ALL{c}\ALL{i}\;\mathrm{Car}(c) \land \mathrm{AtIntersection}(c,i) \land \mathrm{GreenLight}(i) \land \neg\mathrm{PedestrianCrossing}(i) \ra \mathrm{Proceeds}(c,i)\]

\textbf{(iii)} Stop if obstacle on road:
\[\ALL{c}\;\mathrm{Car}(c) \land \mathrm{ObstacleOnRoad}(c) \ra \mathrm{Stops}(c)\]

\textbf{(iv)} Reduce speed in school zone during active hours:
\[\ALL{c}\ALL{z}\;\mathrm{Car}(c) \land \mathrm{InSchoolZone}(c,z) \land \mathrm{ActiveHours}(z) \ra \mathrm{ReduceSpeed}(c)\]

\textbf{(v)} Change lanes only if adjacent lane clear and car is moving:
\[\ALL{c}\ALL{l}\;\mathrm{Car}(c) \land \mathrm{AdjacentLaneClear}(c,l) \land \mathrm{Moving}(c) \ra \mathrm{ChangeLane}(c,l)\]
\end{solbox}

\subsection{Part (b): Drone Delivery FOL + Forward Chaining (4 marks)}

\begin{solbox}{Drone FOL + Forward Chaining --- 2023 Q4(b)}
\textbf{Predicate declarations:}
\begin{itemize}[noitemsep]
  \item $\mathrm{At}(d,l)$: drone $d$ is at location $l$
  \item $\mathrm{Route}(l_1,l_2)$: there exists an air route from $l_1$ to $l_2$
  \item $\mathrm{ObstacleOnRoute}(l_1,l_2)$: there is an obstacle on the route from $l_1$ to $l_2$
  \item $\mathrm{CanFly}(d,l_1,l_2)$: drone $d$ can fly from $l_1$ to $l_2$
\end{itemize}

\textbf{Knowledge Base (Facts):}
\begin{align*}
  KB_1&: \mathrm{At(drone,\;L_0)} \\
  KB_2&: \mathrm{Route}(L_0,\;L_1) \quad\text{(route R1 --- clear)}\\
  KB_3&: \mathrm{Route}(L_1,\;L_2) \quad\text{(route R2 --- has obstacle!)}\\
  KB_4&: \mathrm{Route}(L_2,\;L_p) \quad\text{(route R3 --- clear)}\\
  KB_5&: \mathrm{ObstacleOnRoute}(L_1,\;L_2)\quad\text{(obstacle on R2)}\\
  KB_6&: \neg\mathrm{ObstacleOnRoute}(L_0,\;L_1)\quad\text{(R1 is clear)}\\
  KB_7&: \neg\mathrm{ObstacleOnRoute}(L_2,\;L_p)\quad\text{(R3 is clear)}
\end{align*}

\textbf{Movement Rule in FOL:}
\[\ALL{d}\ALL{l_1}\ALL{l_2}\;\mathrm{At}(d,l_1) \land \mathrm{Route}(l_1,l_2) \land \neg\mathrm{ObstacleOnRoute}(l_1,l_2) \ra \mathrm{At}(d,l_2)\]

\textbf{Forward Chaining (data-driven):}
\begin{center}
\begin{tabular}{@{}l@{\quad}l@{\quad}l@{}}
  \toprule
  \textbf{Iteration} & \textbf{Antecedents matched} & \textbf{Conclusion (added to KB)} \\
  \midrule
  1 & $\mathrm{At(drone,L_0)} + \mathrm{Route}(L_0,L_1) + \neg\mathrm{Obstacle}(L_0,L_1)$ & $\mathrm{At(drone,L_1)}$ \\[3pt]
  2 & $\mathrm{At(drone,L_1)} + \mathrm{Route}(L_1,L_2) + \mathrm{ObstacleOnRoute}(L_1,L_2)$ & \textbf{RULE FAILS} \\
    & Condition $\neg\mathrm{ObstacleOnRoute}(L_1,L_2)$ is FALSE & No inference \\
  \bottomrule
\end{tabular}
\end{center}

\textbf{Conclusion:} Drone reaches $L_1$ successfully, but \textbf{cannot proceed to $L_2$} due to the obstacle on route $R_2$. Therefore, the drone \textbf{cannot reach the destination $L_p$} via the given routes.

\medskip
\textit{(In a full system, an alternative route planner would re-route via a different path if one existed.)}
\end{solbox}

\newpage
%% ==========================================
\section{2024 --- Q4: ER + Perfect Square FOL (9 marks)}
%% ==========================================

\begin{qbox}{Question --- 2024 Q4}
\begin{enumerate}[label=(\alph*),noitemsep]
  \item In ER (Evidential Reasoning), how is a high-level attribute assessed through associated lower-level attributes? \hfill\textit{(1 mark)}
  \item What is the combined degree of belief in the original ER approach? \hfill\textit{(1 mark)}
  \item Consider the following sentences: \hfill\textit{(7 marks)}
    \begin{enumerate}[label=(\roman*),noitemsep]
      \item If a perfect square is divisible by a prime $P$, then it is also divisible by square of $P$
      \item Every perfect square is divisible by some prime
      \item 36 is a perfect square
      \item Does there exist a prime $q$ such that square of $q$ divides 36?
    \end{enumerate}
    \textbf{A)} Translate the sentences above into FOL sentences. \hfill\textit{(2 marks)}\\
    \textbf{B)} Prove (iv) by using Universal Elimination, Existential Elimination, Modus Ponens, and Existential Introduction. \hfill\textit{(1.5 marks)}\\
    \textbf{C)} Convert the FOL sentences into Clausal Form by using Canonical Form (CF). \hfill\textit{(1.5 marks)}\\
    \textbf{D)} Apply the Resolution method to prove (iv). \hfill\textit{(2 marks)}
\end{enumerate}
\end{qbox}

\subsection{Part (a): Evidential Reasoning --- Assessment of High-Level Attributes (1 mark)}

\begin{solbox}{Evidential Reasoning --- 2024 Q4(a)}
In \textbf{Evidential Reasoning (ER)}, a high-level attribute (e.g., ``Disease D'') is assessed by combining probabilistic evidence from multiple lower-level attributes (e.g., symptoms, test results) using the Dempster--Shafer theory of belief functions.

Each lower-level attribute provides a \textbf{basic probability assignment} (mass) distributed over hypotheses. The masses from all evidence sources are combined using \textbf{Dempster's combination rule} (orthogonal sum), yielding a \textbf{belief measure} $\mathrm{Bel}(H)$ and a \textbf{plausibility measure} $\mathrm{Pl}(H)$ for the high-level hypothesis $H$.
\end{solbox}

\subsection{Part (b): Combined Degree of Belief (1 mark)}

\begin{solbox}{Combined Degree of Belief --- 2024 Q4(b)}
The \textbf{combined degree of belief} in the original ER approach uses \textbf{Dempster's rule}:
\[m_{1\oplus 2}(A) = \frac{1}{1-K}\sum_{B\cap C = A} m_1(B)\cdot m_2(C)\]
where $K = \sum_{B\cap C = \emptyset} m_1(B)\cdot m_2(C)$ is the \textbf{conflict factor} (measure of contradictory evidence).

The term $\frac{1}{1-K}$ normalizes the combined belief to sum to 1.\\
\textbf{Belief} $\mathrm{Bel}(H) = \sum_{A \subseteq H} m(A)$ represents total committed belief; \textbf{Plausibility} $\mathrm{Pl}(H) = 1 - \mathrm{Bel}(\neg H)$ represents the upper bound of belief.
\end{solbox}

\subsection{Part (c-A): FOL Translation (2 marks)}

\begin{solbox}{FOL Translation --- 2024 Q4(c-A)}
\textit{Identical to 2021 Q5(b-i). See Section 3.}

\begin{align*}
  F_1&: \ALL{x}\ALL{p}\;\mathrm{PerfSq}(x) \land \mathrm{Prime}(p) \land \mathrm{Divides}(p,x) \ra \mathrm{Divides}(\mathrm{square}(p),x) \\
  F_2&: \ALL{x}\;\mathrm{PerfSq}(x) \ra \EX{p}\;\mathrm{Prime}(p) \land \mathrm{Divides}(p,x)\\
  F_3&: \mathrm{PerfSq}(36)\\
  \text{Goal}&: \EX{q}\;\mathrm{Prime}(q) \land \mathrm{Divides}(\mathrm{square}(q),36)
\end{align*}
\end{solbox}

\subsection{Part (c-B): Proof Using Inference Rules (1.5 marks)}

\begin{solbox}{Inference Rule Proof --- 2024 Q4(c-B)}
\textbf{Prove:} $\EX{q}\;\mathrm{Prime}(q)\land\mathrm{Divides}(\mathrm{square}(q),36)$

\begin{enumerate}[noitemsep,label=\textbf{\arabic*.}]
  \item $\mathrm{PerfSq}(36)$ \hfill $[F_3]$
  \item $\mathrm{PerfSq}(36) \ra \EX{p}\;\mathrm{Prime}(p)\land\mathrm{Divides}(p,36)$ \hfill $[\mathbf{UE}\;x/36\;\text{on}\;F_2]$
  \item $\EX{p}\;\mathrm{Prime}(p)\land\mathrm{Divides}(p,36)$ \hfill $[\mathbf{Modus\;Ponens:}\;1+2]$
  \item $\mathrm{Prime}(p_0)\land\mathrm{Divides}(p_0,36)$ \hfill $[\mathbf{EE}\;\text{(Skolem const }p_0\text{): from 3}]$
  \item $\mathrm{Prime}(p_0)$ and $\mathrm{Divides}(p_0,36)$ \hfill $[\text{And-Elim from 4}]$
  \item $\mathrm{PerfSq}(36)\land\mathrm{Prime}(p_0)\land\mathrm{Divides}(p_0,36) \ra \mathrm{Divides}(\mathrm{square}(p_0),36)$\hfill $[\mathbf{UE}\;x/36,p/p_0\;\text{on}\;F_1]$
  \item $\mathrm{Divides}(\mathrm{square}(p_0),36)$ \hfill $[\textbf{MP:}\;1,5,6]$
  \item $\mathrm{Prime}(p_0)\land\mathrm{Divides}(\mathrm{square}(p_0),36)$ \hfill $[\text{And-Intro: 5+7}]$
  \item $\EX{q}\;\mathrm{Prime}(q)\land\mathrm{Divides}(\mathrm{square}(q),36)$ \hfill $[\mathbf{EI:}\;\text{from 8, generalise }p_0\;\text{to }q]$
\end{enumerate}
\textbf{Goal proved.} \checkmark
\end{solbox}

\subsection{Part (c-C): Canonical Form Conversion (1.5 marks)}

\begin{solbox}{Clause Form --- 2024 Q4(c-C)}
\textit{Identical to 2021 Q5(b-ii). See Section 3.}

\medskip
\begin{tabular}{@{}l@{\quad}l@{}}
  \toprule
  \textbf{Clause} & \textbf{Source}\\
  \midrule
  $C_1: \cl{\neg\mathrm{PerfSq}(x),\;\neg\mathrm{Prime}(p),\;\neg\mathrm{Divides}(p,x),\;\mathrm{Divides}(\mathrm{square}(p),x)}$ & $F_1$ (Steps 2,7)\\[3pt]
  $C_{2a}: \cl{\neg\mathrm{PerfSq}(x),\;\mathrm{Prime}(g(x))}$ & $F_2$ (Skolem fn $g$)\\[3pt]
  $C_{2b}: \cl{\neg\mathrm{PerfSq}(x),\;\mathrm{Divides}(g(x),x)}$ & $F_2$ (Skolem fn $g$)\\[3pt]
  $C_3: \cl{\mathrm{PerfSq}(36)}$ & $F_3$ \\[3pt]
  $C_4: \cl{\neg\mathrm{Prime}(q),\;\neg\mathrm{Divides}(\mathrm{square}(q),36)}$ & Negated goal \\
  \bottomrule
\end{tabular}
\end{solbox}

\subsection{Part (c-D): Resolution Proof (2 marks)}

\begin{solbox}{Resolution Proof --- 2024 Q4(c-D)}
\textit{Identical to 2021 Q5(b-iii). See Section 3 for full 7-step derivation.}

\medskip
\textbf{Summary:}
\begin{enumerate}[noitemsep,label=\textbf{\arabic*.}]
  \item $C_3 + C_{2a}$ $\{x/36\}$ $\to$ $\mathrm{Prime}(g(36))$
  \item $C_3 + C_{2b}$ $\{x/36\}$ $\to$ $\mathrm{Divides}(g(36),36)$
  \item $C_3 + C_1$ $\{x/36\}$ $\to$ $\neg\mathrm{Prime}(p)\lor\neg\mathrm{Divides}(p,36)\lor\mathrm{Divides}(\mathrm{square}(p),36)$
  \item (1)+(3) $\{p/g(36)\}$ $\to$ $\neg\mathrm{Divides}(g(36),36)\lor\mathrm{Divides}(\mathrm{square}(g(36)),36)$
  \item (2)+(4) $\to$ $\mathrm{Divides}(\mathrm{square}(g(36)),36)$
  \item (1)+$C_4$ $\{q/g(36)\}$ $\to$ $\neg\mathrm{Divides}(\mathrm{square}(g(36)),36)$
  \item (5)+(6) $\to$ $\boxed{\square}$ \textbf{(Contradiction $\Rightarrow$ Goal Proved!)} \checkmark
\end{enumerate}
\end{solbox}

\newpage
%% ==========================================
\section*{Quick Exam Summary --- Exam Patterns for This Topic}
%% ==========================================

\begin{warnbox}{Every Year Patterns --- Memorise These}
\begin{enumerate}[noitemsep]
  \item \textbf{Marcus/Pompeii} (2020, 2022): Always 4 parts --- (a) WFF translation (b) backward reasoning (c) clause form (d) resolution. Now changes year each paper (2021 in 2020 paper, 2023 in 2022 paper).
  \item \textbf{Perfect Square FOL} (2021, 2024): Always same 4 sentences. 2021 asks for UE/EE/MP/EI proof. 2024 adds ER questions. Proof is always: Skolem $g(36)$ → derive Prime$+$Divides → contradict negated goal.
  \item \textbf{FOL Translation: Robot/Drone/Car} (2023): Declare predicates FIRST, then encode each action rule as $\forall$ sentence with $\rightarrow$.
  \item \textbf{KB Resolution (Wumpus)} (2021 Q4d): Use $R_4(\neg B_{1,1})$ + clause from biconditional of $R_2$ to immediately derive contradiction. 2-step proof.
  \item \textbf{Satisfiability check} (2021 Q4c, 2022 A-4a): Simplify algebraically first (use Distributivity, De Morgan). Then check truth table.
  \item \textbf{Clause form algorithm asked directly} (2022, 2016): Know all 9 steps. Pay attention to Skolemization rule (constant vs.\ function depends on outer $\forall$).
\end{enumerate}
\end{warnbox}

\end{document}
